% !TEX root = ../../../main.tex
% 来源：中学数学实验教材·第二册（几何）
% 章节：实验几何 / 相似和相似比
% 原始文件：1.tex，行 3131-3168
% 类型：tikz
% ---
\begin{tikzpicture}[scale=.5]
	\begin{scope}
\tkzDefPoints{0/0/A,3/0/B, 2/2/C}
\tkzDrawPolygon(A,B,C)
\tkzLabelPoints[below](A,B)
\tkzLabelPoints[above](C)
\end{scope}
\begin{scope}[xshift=4.5cm]
\tkzDefPoints{0/0/A',9/0/B', 6/6/C'}
\tkzDrawPolygon(A',B',C')
\tkzLabelPoints[below](A',B')
\tkzLabelPoints[above](C')
\tkzDefPointWith[linear, K=.333](A',B') \tkzGetPoint{M_1}
\tkzDefPointWith[linear, K=.667](A',B') \tkzGetPoint{M_2}

\tkzDefPointWith[linear, K=.333](A',C') \tkzGetPoint{M_1'}
\tkzDefPointWith[linear, K=.667](A',C') \tkzGetPoint{M_2'}

\tkzDefPointWith[linear, K=.333](C',B') \tkzGetPoint{M_1''}
\tkzDefPointWith[linear, K=.667](C',B') \tkzGetPoint{M_2''}
\tkzDrawSegments(M_2'',M_1' M_1'',M_2' M_1,M_1' M_2,M_2' M_1'',M_1 M_2'',M_2)
\tkzInterLL(M_1'',M_1)(M_2,M_2') \tkzGetPoint{O}
\tkzLabelPoints[below](M_1,M_2,O)
\tkzLabelPoints[left](M_1',M_2')
\tkzLabelPoints[right](M_1'',M_2'')
\end{scope}
\begin{scope}[xshift=15cm]
\tkzDefPoints{0/0/A'',3/0/B, 2/2/C}
\tkzDefPointWith[linear, K=2](A'',B) \tkzGetPoint{B''}
\tkzDefPointWith[linear, K=2](A'',C) \tkzGetPoint{C''}
\tkzDrawPolygon(A'',B'',C'')
\tkzDefMidPoint(C'',B'') \tkzGetPoint{D}
\tkzDrawPolygon(D,B,C)
\tkzLabelPoints[above](C'')
\tkzLabelPoints[below](A'', B'')

\end{scope}
\end{tikzpicture}
